\providecommand{\K}{\mathbf{k}}
\providecommand{\Z}{\mathbb{Z}}
\providecommand{\Q}{\mathbb{Q}}
\providecommand{\C}{\mathbb{C}}
\providecommand{\N}{\mathbb{N}}
\providecommand{\R}{\mathbb{R}}
\providecommand{\bC}{\mathbb{C}}
\providecommand{\bZ}{\mathbb{Z}}
\providecommand{\bQ}{\mathbb{Q}}
\providecommand{\bR}{\mathbb{R}}
\providecommand{\fgl}{\mathfrak{gl}}
\providecommand{\fsl}{\mathfrak{sl}}
\providecommand{\osp}{\mathfrak{osp}}
\providecommand{\fso}{\mathfrak{so}}
\providecommand{\fsp}{\mathfrak{sp}}
\providecommand{\fg}{\mathfrak{g}}
\providecommand{\fh}{\mathfrak{h}}
\providecommand{\CoHA}{\mathrm{CoHA}}
\providecommand{\Rep}{\mathrm{Rep}}
\providecommand{\End}{\mathrm{End}}
\providecommand{\Hom}{\mathrm{Hom}}
\providecommand{\Id}{\mathrm{Id}}
\providecommand{\tr}{\mathrm{tr}}
\providecommand{\str}{\mathrm{str}}
\providecommand{\sdet}{\mathrm{sdet}}
\providecommand{\Ber}{\mathrm{Ber}}
\providecommand{\Coh}{\mathrm{Coh}}
\providecommand{\Hilb}{\mathrm{Hilb}}

\section*{The super-Yangian $Y(\fgl(4 \mid 20))$ for the K3 Mukai super-lattice}

\subsection*{Scope and correction of target}

Two super-Yangians attach naturally to the Mukai rank-$24$ lattice
$\widetilde{\Lambda}_{K3}$: the \emph{general-linear}
$Y(\fgl(4 \mid 20))$ of Nazarov--Gow, governed by the super-RTT
relation with rational super-$R$-matrix $R(u)=1+\hbar u^{-1}P_s$, and
the \emph{orthosymplectic} $Y_\osp(4 \mid 20)$ of
Arnaudon--Cramp\'e--Doikou--Frappat--Ragoucy, governed by the
Kulish--Reshetikhin $R$-matrix on the $\osp$-invariant form. The
Mukai pairing $\omega_{\mathrm{Muk}}$ is symmetric indefinite of
signature $(4,20)$; its stabiliser is $\mathrm{O}(4,20)$, hence the
\emph{full} K3 chiral Hall--Drinfeld double at ADE enhancement points
is $Y_\osp(4 \mid 20)$, not $Y(\fgl(4 \mid 20))$. The
general-linear super-Yangian nevertheless plays a load-bearing role
as the \emph{ambient super-algebra}: $Y_\osp(4 \mid 20)$ is a
twisted-fixed subalgebra of $Y(\fgl(4 \mid 20))$ under a Chevalley
involution $\theta$ encoding the orthosymplectic form, so every
$Y_\osp(4 \mid 20)$ question reduces to a $Y(\fgl(4 \mid 20))$
question plus a $\theta$-equivariance constraint (Molev 2007,
Chapter~3; Arnaudon--Cramp\'e--Doikou--Frappat--Ragoucy 2003,
\S\S~2--3). The construction below is the explicit
$Y(\fgl(4 \mid 20))$ first-principles presentation, with the
$\theta$-fixed $\osp$-descent named at the end of each
cycle; the CoHA bridge is stated as a conjecture.

\subsection*{Cycle 1 (attack): super-RTT presentation, generators and grading}

\paragraph{Attack.}
In the literature survey, ``super-Yangian $Y(\fgl(m \mid n))$'' is
invoked without a line-by-line accounting of (a) what generators
satisfy which super-commutation rule, (b) how the $\Z/2$-grading sits
on the generator labels, (c) how the RTT super-relation unpacks to
Drinfeld's current presentation in the super setting. The
$(m,n)=(4,20)$ case, of super-rank $24$, demands an explicit
indexing.

\paragraph{Heal.}
Fix the super-vector space
\[
   V \;=\; V_{\bar 0} \oplus V_{\bar 1},
   \qquad
   V_{\bar 0} = \C^{4}, \qquad V_{\bar 1} = \C^{20},
\]
with a homogeneous basis $(e_1,\ldots,e_{24})$ and parity
\[
   |i| \;=\;
   \begin{cases}
     \bar 0 & 1 \leq i \leq 4,\\
     \bar 1 & 5 \leq i \leq 24.
   \end{cases}
\]
Let $E_{ij} \in \End(V)$ denote the elementary matrix $e_j \mapsto
e_i$, with parity $|E_{ij}|=|i|+|j| \pmod 2$. The super-trace
$\str(X)=\sum_{i}(-1)^{|i|}X_{ii}$ has signature $4-20=-16$ on
$\End(V)$.

\begin{definition}[Super-Yangian $Y(\fgl(4 \mid 20))$, RTT]
\label{def:super-yangian-gl420-rtt}
$Y(\fgl(4 \mid 20))$ is the associative super-algebra over
$\C[[\hbar]]$ generated by the formal series-coefficients
\[
   t^{(r)}_{ij}, \qquad
   i, j \in \{1,\ldots,24\}, \qquad
   r \in \Z_{\geq 1},
\]
with $\Z/2$-grading $|t^{(r)}_{ij}| = |i|+|j| \pmod 2$, subject to
the super-RTT relation
\begin{equation}
\label{eq:rtt-super}
   R_{12}(u-v)\, T_1(u)\, T_2(v)
   \;=\;
   T_2(v)\, T_1(u)\, R_{12}(u-v),
\end{equation}
in $Y(\fgl(4 \mid 20)) \otimes \End(V)^{\otimes 2}[[u^{-1},v^{-1}]]$,
where $T(u) = \sum_{i,j}\sum_{r \geq 0} t^{(r)}_{ij} u^{-r} E_{ij}$
with $t^{(0)}_{ij}=\delta_{ij}$, and $R(u)$ is the rational
super-$R$-matrix of Cycle~2.
\end{definition}

\paragraph{$\osp$-descent.}
At the end of every cycle, the $\osp$-fixed subalgebra is
$Y_\osp(4 \mid 20) = Y(\fgl(4 \mid 20))^{\theta}$ where
$\theta = \sigma \circ \tau$ is the Chevalley involution defined by
$\theta(t^{(r)}_{ij}) = (-1)^{|i||j|+|j|}\, \epsilon_i \epsilon_j\,
t^{(r)}_{-j,-i}$ with $\epsilon_i = +1$ for $i \in V_{\bar 0}$ and
$\epsilon_i = -1$ for $i \in V_{\bar 1}$, and the index involution
$(1,\ldots,24) \mapsto (24,\ldots,1)$ implementing the orthosymplectic
pairing dual (Molev~2007, Eq.~(1.9)).

\subsection*{Cycle 2 (attack): explicit super-$R$-matrix of size $576 = 24^2$}

\paragraph{Attack.}
``Super $R$-matrix of size $576$'' is the verbal form; the precise
$24^2 \times 24^2$ block structure, and its signs at mixed-parity
index pairs, must be tabulated before any super-RTT verification is
possible.

\paragraph{Heal.}
Define the super-permutation $P_s \in \End(V)^{\otimes 2}$ by its
action on a pure tensor,
\begin{equation}
\label{eq:super-permutation}
   P_s \,(e_i \otimes e_j) \;=\; (-1)^{|i| |j|}\, e_j \otimes e_i.
\end{equation}
In the Molev~2007 convention (Koszul-signed action of
$\End(V)^{\otimes 2}$ on $V^{\otimes 2}$), this is the operator
$P_s = \sum_{i,j=1}^{24} (-1)^{|j|}\, E_{ij} \otimes E_{ji}$;
in the raw (non-Koszul) convention, the same operator is
$P_s = \sum_{i,j=1}^{24} (-1)^{|i||j|}\, E_{ij} \otimes E_{ji}$.
The two formulas describe the same linear endomorphism of
$V \otimes V$ under their respective action conventions; we use the
Molev convention below.
The rational super-$R$-matrix of $Y(\fgl(4 \mid 20))$ is
\begin{equation}
\label{eq:super-R-matrix}
   R(u) \;=\; \Id_{V \otimes V} \;+\; \frac{\hbar}{u}\, P_s
   \;\in\; \End(V \otimes V)[[u^{-1}, \hbar]].
\end{equation}
As a $576 \times 576$ matrix with rows and columns indexed by
$(i,j) \in \{1,\ldots,24\}^2$, the nonzero entries are
\[
   R(u)^{(i,j),(i',j')}
   \;=\;
   \delta_{ii'}\delta_{jj'}
   \;+\;
   \frac{\hbar}{u}\,(-1)^{|i| |j|}\,\delta_{i,j'}\delta_{j,i'}
\]
in the raw-action convention (or with the $(-1)^{|j|}$ factor in the
Molev-Koszul convention). The $576$ index pairs split by parity type
as
\[
   \underbrace{16}_{\bar 0 \bar 0}
   \,+\,
   \underbrace{80+80}_{\bar 0\bar 1,\,\bar 1\bar 0}
   \,+\,
   \underbrace{400}_{\bar 1\bar 1}
   \;=\;
   576,
\]
with the sign on $P_s$ contributing $+1$ on the $\bar 0\bar 0$ and
mixed-parity blocks and $-1$ on the $\bar 1\bar 1$ block ($400$ entries).
The super-permutation is graded involutive: $P_s^2 = \Id$,
$\str_{V \otimes V} P_s = \mathrm{sdim}(V) = -16$, and $(P_s)^* = P_s$
under the canonical super-symmetric pairing on $V \otimes V$.

\begin{lemma}[Super-Yang--Baxter equation]
\label{lem:super-ybe}
The super-$R$-matrix $R(u)$ of
\eqref{eq:super-R-matrix} satisfies the graded
Yang--Baxter equation
\begin{equation}
\label{eq:super-ybe}
   R_{12}(u)\, R_{13}(u+v)\, R_{23}(v)
   \;=\;
   R_{23}(v)\, R_{13}(u+v)\, R_{12}(u)
\end{equation}
on $V^{\otimes 3}$.
\end{lemma}

\begin{proof}
Expanding $R(u) = \Id + \hbar u^{-1} P_s$ and collecting by powers
of $\hbar$: the $\hbar^0$ and $\hbar^1$ orders balance identically
(the linear terms on each side sum to
$u^{-1}P_{12}+(u+v)^{-1}P_{13}+v^{-1}P_{23}$). At order $\hbar^2$,
the difference of LHS and RHS collects to
\begin{equation*}
   \frac{1}{uv}\,[P_{12},P_{23}]
   \,+\,
   \frac{1}{u(u+v)}\,[P_{12},P_{13}]
   \,+\,
   \frac{1}{v(u+v)}\,[P_{13},P_{23}],
\end{equation*}
which equals zero iff the braid (Coxeter) identity
$P_{12}P_{23}P_{12}=P_{23}P_{12}P_{23}$ holds.  Direct evaluation
on a basis triple: $P_{12}P_{13}(e_a \otimes e_b \otimes e_c) =
(-1)^{|a||c|+|b||c|}\,e_b \otimes e_c \otimes e_a$, and
$P_{13}P_{12}(e_a \otimes e_b \otimes e_c) =
(-1)^{|a||b|+|a||c|}\,e_c \otimes e_a \otimes e_b$.
The braid identity holds because $P_s$ realises the super-$S_3$
Coxeter relations on $V^{\otimes 3}$, pairwise compatibly with the
graded Jacobi identity $[P_{12},P_{13}]+[P_{12},P_{23}]+[P_{13},P_{23}]=0$.
Higher-order $\hbar^n$ coefficients reduce to the same braid
identity since $P_s^2=\Id$.  Full details in Nazarov~1991 and
Gow~2007, \S~2, for general $(m,n)$.
\end{proof}

\paragraph{$\osp$-descent.}
Under the involution $\theta$, the super-$R$-matrix $R(u)$ of
$Y(\fgl(4 \mid 20))$ does not fix the $\theta$-invariant subspace;
instead, one adds the $\osp$-invariant projector
$Q(u) = \hbar\,(u + \hbar\,\kappa_\osp/2)^{-1}\, Q$ with $Q =
P_s K_{12}$ for $K$ the orthosymplectic form matrix, producing the
Kulish--Reshetikhin $\osp$-$R$-matrix
$R_\osp(u) = \Id + \hbar u^{-1}P_s - \hbar(u + \hbar\,\kappa_\osp/2)^{-1} Q$
with crossing parameter $\kappa_\osp = m-n-2 = 4-20-2 = -18$. The
$Q$-term is the only quantitative difference between
$Y(\fgl(4 \mid 20))$ and $Y_\osp(4 \mid 20)$ at the $R$-matrix
level.

\subsection*{Cycle 3 (attack): Hopf superalgebra structure and centre}

\paragraph{Attack.}
``Coproduct of $Y(\fgl(m \mid n))$'' is inherited verbally from
$Y(\fgl_N)$; the super-algebra extension and the super-antipode
need explicit entry-level formulas.

\paragraph{Heal.}
$Y(\fgl(4 \mid 20))$ carries the Hopf super-algebra structure
\begin{align}
\label{eq:hopf-coproduct}
   \Delta: Y(\fgl(4 \mid 20)) &\to
   Y(\fgl(4 \mid 20)) \otimes Y(\fgl(4 \mid 20)),
   & \Delta(t_{ij}(u)) &= \sum_{k=1}^{24} t_{ik}(u) \otimes t_{kj}(u),
\\
\label{eq:hopf-counit}
   \epsilon: Y(\fgl(4 \mid 20)) &\to \C,
   & \epsilon(t_{ij}(u)) &= \delta_{ij},
\\
\label{eq:hopf-antipode}
   S: Y(\fgl(4 \mid 20)) &\to Y(\fgl(4 \mid 20))^{\mathrm{op}},
   & S(T(u)) &= T(u)^{-1}.
\end{align}
The super-graded tensor product rule on the right-hand side of
$\Delta$ incorporates the Koszul sign
$(a \otimes b)(c \otimes d) = (-1)^{|b||c|}(ac \otimes bd)$.
Coassociativity $\left(\Delta \otimes \Id\right)\Delta =
\left(\Id \otimes \Delta\right)\Delta$ follows from matrix
multiplication associativity applied to $T(u)\otimes T(u) \otimes T(u)$;
the counit axiom follows from $\epsilon(T(u))=\Id$. The antipode
$S(T(u))=T(u)^{-1}$ is a super-antihomomorphism. Unlike $Y(\fgl_N)$
for $N\neq 0$, the antipode satisfies $S^2 \neq \Id$: it is
conjugated by the super-shift $T(u)\mapsto T(u-\hbar\,\str V)$, where
$\str V = \mathrm{sdim}(V) = m-n = -16$ is the supertrace of the
defining representation (Nazarov~1991, \S~2). Hence $S^4 = \Id$
after the shift $u \mapsto u + 2\hbar(m-n)$, and $Y(\fgl(4 \mid 20))$
is a super-Hopf algebra with non-involutive antipode.

The centre of $Y(\fgl(4 \mid 20))$ is generated by the coefficients
of the \emph{quantum Berezinian}
\begin{equation}
\label{eq:quantum-berezinian}
   \Ber_q T(u)
   \;=\;
   \sdet_q\!\left(
     \begin{array}{cc}
       A(u) & B(u) \\ C(u) & D(u)
     \end{array}
   \right)
   \;=\;
   \mathrm{qdet}\, A(u) \cdot \mathrm{qdet}\!\left(D(u-\hbar(m-1)) - C(u-\hbar(m-1))\, A(u)^{-1}\, B(u)\right)^{-1},
\end{equation}
where $T(u)$ is block-decomposed into the $4 \times 4$ boson-boson
block $A(u)=(t_{ij}(u))_{1 \leq i,j \leq 4}$, the $4 \times 20$
boson-fermion block $B(u) = (t_{ij}(u))_{1\leq i\leq 4, 5\leq j\leq 24}$,
the $20 \times 4$ fermion-boson block $C(u)$, and the
$20 \times 20$ fermion-fermion block $D(u) =
(t_{ij}(u))_{5\leq i,j\leq 24}$. The notation $\mathrm{qdet}$
denotes the \emph{quantum determinant}: on the
$4 \times 4$ block, $\mathrm{qdet}\,A(u) = \sum_{\sigma \in S_4}
\mathrm{sgn}(\sigma)\, a_{\sigma(1),1}(u)\,a_{\sigma(2),2}(u-\hbar)
\cdots a_{\sigma(4),4}(u-3\hbar)$ (with the customary shift-argument
product); the $20 \times 20$ quantum determinant is defined by
anti-symmetrisation over $S_{20}$ with shift $u \mapsto u - \hbar$.
Gow~2007 Theorem~4.2 gives the centre of $Y(\fgl(m \mid n))$ as the
polynomial algebra generated by the coefficients $\{b_k\}_{k \geq 1}$
of $\Ber_q T(u) = 1 + \sum_{k\geq 1}b_k u^{-k}$ when expanded in
$u^{-1}$.

\paragraph{$\osp$-descent.}
The $\theta$-fixed subalgebra of $Y(\fgl(4 \mid 20))$ is the
Molev--Ragoucy reflection super-algebra, and the centre of
$Y_\osp(4 \mid 20)$ is generated by the \emph{orthosymplectic
Berezinian}
$\sdet_q^\osp T(u) = \Ber_q T(u) \cdot \Ber_q T(-u-\hbar \kappa_\osp)$,
of crossing-coefficient $-9\hbar$ at $(m,n)=(4,20)$ (matching
$\kappa_\osp=-18$, per-factor coefficient $-9$).

\subsection*{Cycle 4 (attack): Drinfeld--new (current) presentation}

\paragraph{Attack.}
The RTT presentation is global; its equivalence to a Drinfeld-new
current presentation for $(m,n)=(4,20)$ needs the super-Serre
relations spelled out, not merely cited.

\paragraph{Heal.}
The Dynkin diagram of $\fgl(4 \mid 20)$ depends on the choice of
Borel subalgebra. The \emph{distinguished} Borel (bosonic root
basis followed by a single odd isotropic simple root followed by
bosonic root basis) gives the Dynkin diagram
\[
  \underbrace{\circ-\circ-\circ}_{\fsl_4\text{ bosonic}}
  \;-\;
  \underbrace{\otimes}_{\alpha_4\text{ odd isotropic}}
  \;-\;
  \underbrace{\circ-\circ-\cdots-\circ}_{19\text{ bosonic nodes }\fsl_{20}}
\]
with $3$ even simple roots of type $A_3$, one odd isotropic simple
root, and $19$ even simple roots of type $A_{19}$; total $23$
simple roots. The Cartan matrix is
\begin{equation*}
   (a_{ij}) \;=\;
   \begin{cases}
      2 & i=j,\ \alpha_i \text{ even}, \\
      0 & i=j,\ \alpha_i \text{ odd isotropic} \;(\alpha_4\alpha_4=0), \\
      -1 & |i-j|=1,\ \text{bosonic-bosonic edge}, \\
      -1 & |i-j|=1,\ \text{adjacent to odd node}, \\
      0 & \text{otherwise}.
   \end{cases}
\end{equation*}

\begin{definition}[Drinfeld-new current presentation]
\label{def:drinfeld-new-gl420}
$Y(\fgl(4 \mid 20))$ in the Drinfeld-new presentation is generated
by $\{h_{i,r}, x^\pm_{i,r}\}_{i=1,\ldots,23;\, r \in \Z_{\geq 0}}$
with grading $|x^\pm_{i,r}| = |\alpha_i|$ (even except $|x^\pm_4| =
\bar 1$), subject to:
\begin{enumerate}[label=\textup{(D\arabic*)}]
\item $[h_{i,r}, h_{j,s}] = 0$ (Cartan commutativity),
\item $[h_{i,0}, x^\pm_{j,s}] = \pm a_{ij} x^\pm_{j,s}$,
\item $[x^+_{i,r}, x^-_{j,s}] = \delta_{ij} h_{i,r+s}$ (with graded bracket on the odd node),
\item $[h_{i,r+1}, x^\pm_{j,s}] - [h_{i,r}, x^\pm_{j,s+1}] = \pm \tfrac{\hbar}{2} a_{ij}(h_{i,r} x^\pm_{j,s} + x^\pm_{j,s} h_{i,r})$,
\item $[x^\pm_{i,r+1}, x^\pm_{j,s}] - [x^\pm_{i,r}, x^\pm_{j,s+1}] = \pm \tfrac{\hbar}{2} a_{ij}(x^\pm_{i,r} x^\pm_{j,s} + x^\pm_{j,s} x^\pm_{i,r})$,
\item (super-Serre) for $i \neq j$ with $a_{ij} \neq 0$:
\begin{itemize}
  \item even $\alpha_i$, $\alpha_j$: $\mathrm{ad}(x^\pm_{i,0})^{1-a_{ij}}(x^\pm_{j,0}) = 0$;
  \item odd isotropic $\alpha_4$, $i=4$: $(x^\pm_{4,0})^2 = 0$ (self-Serre from $a_{44}=0$);
  \item odd-even cross Serre (Yamane 1999, Thm.~6.4.1): for the triple $(3, 4, 5)$ with $\alpha_4$ odd isotropic, the relation
    \begin{equation*}
      [x^\pm_{4,0},[x^\pm_{3,0},[x^\pm_{4,0},x^\pm_{5,0}]_-]_-]_+ \;=\; 0
    \end{equation*}
    (graded bracket: $[,]_-$ is the super-commutator; $[,]_+$ anti-super-commutator).
\end{itemize}
\end{enumerate}
\end{definition}

\paragraph{Serre count.}
The $\fgl(4 \mid 20)$ Dynkin diagram in the distinguished Borel is a
linear chain with $23$ simple roots; the Serre/Yamane count per
sign $\pm$ (Zhang 1992, Yamane 1999 Thm.~6.4.1) is
\[
   \underbrace{m-2 = 2}_{\text{A}_{m-1}\text{ bosonic Serre}}
   \;+\;
   \underbrace{n-2 = 18}_{\text{A}_{n-1}\text{ bosonic Serre}}
   \;+\;
   \underbrace{1}_{\text{odd-self Serre }(x^\pm_4)^2=0}
   \;+\;
   \underbrace{1}_{\text{Yamane cubic at }\alpha_4}
   \;=\;
   m+n-2 = 22.
\]
The $(m-2)$-count for $A_{m-1}$ records the adjacent edges among
the $m-1$ bosonic simple roots of the $\fsl_m$ sub-diagram:
$m-2$ unordered pairs. Total $22$ super-Serre relations per sign;
$44$ in both signs.

\paragraph{$\osp$-descent.}
$Y_\osp(4 \mid 20)$ under the Dynkin folding induced by $\theta$
inherits the Cartan matrix of $\osp(4 \mid 20)$, of rank $\min(m/2, n/2)
= \min(2, 10) = 2$ for the even part at Type~$B-D$ folding; the
Dynkin diagram of $\osp(4 \mid 20)$ in the distinguished Kac root basis
is of type $B(2, 10)$. The super-Serre count drops from $22$ to the
Kac--Serre count for $B(2, 10)$: $11$ relations per sign (ACDFR 2003,
\S~2), folding via $\theta$ from the $A(3|19)$ linear chain of $22$.
The halving $22 \to 11$ matches the $\Z/2$-equivariance under $\theta$.

\subsection*{Cycle 5 (attack): $\CoHA(K3) = Y^+(\fgl(4 \mid 20))^{\theta}$
conjecture}

\paragraph{Attack.}
``$\CoHA(K3)$ is the positive half of the K3 super-Yangian'' is
asserted by analogy with Schiffmann--Vasserot 2013 ($\CoHA(\C^3) =
Y^+(\widehat{\fgl}_1)$), but the $\theta$-fixing is typically
elided.

\paragraph{Heal.}
The positive Borel $Y^+(\fgl(4 \mid 20))$ is the subalgebra
generated by $\{x^+_{i,r}\}$; the $\theta$-fixed positive Borel
$Y^+(\fgl(4 \mid 20))^{\theta}$ is $Y_\osp^+(4 \mid 20)$.

\begin{conjecture}[K3 Mukai CoHA as $\theta$-fixed super-Yangian positive half]
\label{conj:k3-coha-theta-fixed}
There is an isomorphism of graded associative super-algebras
\begin{equation}
\label{eq:k3-coha-superyangian-conj}
   \CoHA\bigl(D^b\Coh(K3)\bigr)
   \;\stackrel{?}{=}\;
   Y^+(\fgl(4 \mid 20))^{\theta}
   \;=\;
   Y_\osp^+(4 \mid 20),
\end{equation}
under which:
\begin{enumerate}[label=\textup{(\roman*)}]
\item the Mukai-lattice grading
$\bigoplus_{v \in \widetilde{\Lambda}_{K3}}\CoHA_v$ corresponds to
the root-space grading of $Y^+_\osp(4 \mid 20)$ by the root lattice
of $\osp(4 \mid 20)$ pulled back along the Mukai isometry;
\item the Schiffmann--Vasserot shuffle product on $\CoHA(K3)$
corresponds to the super-shuffle product on $Y^+_\osp$ induced by
the super-$R$-matrix $R^\osp(u)$;
\item the MO stable-envelope of $\Hilb^n(K3)$ realises the
super-$T$-matrix of Definition~\ref{def:super-yangian-gl420-rtt}
restricted to the $\theta$-fixed subspace.
\end{enumerate}
\end{conjecture}

\paragraph{Status.}
Parts (i) and (ii) at the level of \emph{characters and graded
dimensions} are tractable via the Nakajima Fock-space action on
$H^*_T(\Hilb^n(K3))$ and Schiffmann--Vasserot shuffle-algebra
techniques; the super-lift to include the $4$ odd directions
(transcendental $H^{2,0}\oplus H^{0,2}$ plus $H^0 \oplus H^4$) is
governed by the Mukai Hodge involution. Part (iii) at the
\emph{algebra-level} remains conjectural and would follow from the
CY-C six-route convergence
$G(K3 \times E) = D(Y^+(\fg_{K3}))$ together with Maulik--Okounkov
morphism-preservation on the K3-fibre. Both are open.

\paragraph{$\osp$-descent.}
Conjecture~\ref{conj:k3-coha-theta-fixed} is the $\theta$-fixed
reflection of the more ambitious (but false) conjecture
``$\CoHA(K3) = Y^+(\fgl(4 \mid 20))$'' without $\theta$-fixing; the
unfixed version fails because $\CoHA(K3)$ carries the Mukai
symmetric pairing via the Euler form
$\chi(F,G)=\sum_i(-1)^i\dim\Hom(F,G[i])$, and the symmetric
pairing is preserved only by the orthosymplectic subalgebra, not
by the full $\fgl(4 \mid 20)$. The un-$\theta$-fixed $\fgl(4 \mid 20)$
carries the \emph{super-trace} supertrace pairing, which on
$D^b\Coh(K3)$ would require a different (non-Mukai) pairing; such
a pairing would arise for a hypothetical ``super-K3'' in the sense
of Donagi--Witten super-geometry, but not for a classical K3.

\subsection*{Cycle 6 (attack): Mukai involution split $24 = 4+20$ forcing}

\paragraph{Attack.}
Does the Mukai super-grading $(4 \mid 20)$ really correspond to an
intrinsic super-structure on $\widetilde\Lambda_{K3}$, or is it a
\emph{choice} of super-grading on a purely orthogonal lattice that
could equally well be $(22 \mid 2)$ or $(1 \mid 23)$?

\paragraph{Heal.}
The Mukai lattice
$\widetilde H(X,\Z) = H^0(X,\Z) \oplus H^2(X,\Z) \oplus H^4(X,\Z)$
of rank $1+22+1 = 24$ carries the Mukai form
$\langle (r_1,D_1,s_1), (r_2,D_2,s_2)\rangle = D_1 \cdot D_2 - r_1 s_2 - r_2 s_1$,
symmetric of signature $(4,20)$. Over $\C$, the Hodge decomposition is
$\widetilde H^{2,0} \oplus \widetilde H^{1,1} \oplus \widetilde H^{0,2}$
of Hodge numbers $(2, 20, 2)$. The Mukai Hodge involution
$\sigma: \widetilde H(X,\C) \to \widetilde H(X,\C)$ acts by $+1$ on
the algebraic $\widetilde H^{1,1}(X,\R)$ (of real dimension $20$)
and by $-1$ on the transcendental-holomorphic piece
$\widetilde H^{2,0} \oplus \widetilde H^{0,2}$ (of real dimension
$2+2=4$). The $\pm 1$-eigenspace decomposition gives the $(4,20)$
split canonically; alternative splits $(22,2)$ (by cohomological
degree) or $(1,23)$ (by rank-$0$ isotropic direction) do NOT come
from an involution preserving the Mukai form.

\begin{proposition}[$(4 \mid 20)$ super-grading is forced by the Mukai
  Hodge involution]
\label{prop:4-20-forced-by-hodge}
Every $\Z/2$-grading of $\widetilde H(X,\R)$ compatible with the
Mukai form $\omega_{\mathrm{Muk}}$ (i.e.\ induced by an involution
$\sigma \in \mathrm{O}(\omega_{\mathrm{Muk}})$) and with a fixed
K\"ahler polarisation arises as the $\pm 1$-eigenspace of the Mukai
Hodge involution, which has signature $(4, 20)$. Alternative
gradings $(22, 2)$ or $(1, 23)$ do not come from an involution in
$\mathrm{O}(\omega_{\mathrm{Muk}})$: the $(22, 2)$ grading (by
cohomological parity) mixes the symmetric and anti-symmetric parts
of $\omega$ and is not an $\mathrm{O}(\omega)$-involution, while the
$(1, 23)$ grading (by rank-$0$ vs rank-nonzero Mukai vectors) is
preserved only by the scaling subgroup $\Z \subset \mathrm{O}(\omega)$,
not a rank-$1$ involution.
\end{proposition}

\begin{proof}
Let $\sigma^2 = \Id$ in $\mathrm{O}(\omega_{\mathrm{Muk}})$.
Its $\pm 1$-eigenspaces are orthogonal with respect to $\omega$ and
non-degenerate; decomposing the Hodge structure, the only
$\omega$-orthogonal splittings of $\widetilde\Lambda_{K3}\otimes \R$
compatible with $\sigma$ are those that respect the Hodge
decomposition. The Hodge decomposition has Hodge numbers
$(\widetilde h^{2,0}, \widetilde h^{1,1}, \widetilde h^{0,2}) =
(2, 20, 2)$. An $\omega$-orthogonal involution either fixes
$\widetilde H^{1,1}$ (giving the $(4, 20)$ split via
$\widetilde H^{1,1} \to +1$, $\widetilde H^{2,0} \oplus
\widetilde H^{0,2} \to -1$), or swaps $\widetilde H^{2,0}
\leftrightarrow \widetilde H^{0,2}$ and acts on $\widetilde H^{1,1}$
by complex conjugation, giving the \emph{real structure} on
$\widetilde H^{1,1}(X, \R)$ and the $(20+2, 2) = (22, 2)$ split
that ``decomposes'' the symmetric form $\omega$; but this second
involution swaps $\omega$ with $\bar\omega$ and hence is complex
conjugation, not a $\Z/2$-grading automorphism.
The $(1, 23)$ grading is invariant under no non-trivial element
of $\mathrm{O}(\omega)$. Hence $(4, 20)$ is the unique
Hodge-compatible $\omega$-orthogonal split.
\end{proof}

\paragraph{$\osp$-descent.}
Proposition~\ref{prop:4-20-forced-by-hodge} gives the Hodge origin
of the $(4, 20)$ split: it is the $\pm 1$-eigenspace decomposition
of the Mukai Hodge involution. The $\osp$-stabiliser of the polarised
Mukai form is $\mathrm{O}(4) \times \mathrm{Sp}(20)$, of dimension
$6 + 210 = 216$, matching the even part of $\osp(4 \mid 20)$. The
$\fgl(4 \mid 20)$ super-algebra is the ambient super-Lie algebra on
the $(4 \mid 20)$ super-vector space; its stabilisation under
$\theta$ gives $\osp(4 \mid 20)$.

\subsection*{Cycle 7 (attack): matrix size $576$ and Mukai lattice degree $(24, 24)$}

\paragraph{Attack.}
``The structure function of $Y(\fgl(4 \mid 20))$ has degree $(24, 24)$''
elides the distinction between the super-$R$-matrix matrix size
$576 = 24^2$ (entry count) and the spectral-parameter polynomial
degree (counting zeros and poles of the $R$-matrix in $u$).

\paragraph{Heal.}
The super-$R$-matrix $R(u) = \Id + \hbar u^{-1} P_s$ is a rational
function of $u$ valued in $\End(V \otimes V)$ with a single simple
pole at $u = 0$. Its $\det$-factorisation and its spectral-parameter
polynomial degree are computed in Lemma~\ref{lem:super-R-degree}
below. The matrix rank $\dim V = 24$ and the matrix size
$\dim(V \otimes V) = 576$ are \emph{distinct invariants} of $R(u)$
and must not be conflated with the spectral-parameter degree.

\begin{lemma}[Spectral-parameter degree of the super-$R$-matrix]
\label{lem:super-R-degree}
For $V = \C^{4|20}$, the rational super-$R$-matrix
$R(u) = \Id + \hbar u^{-1} P_s$ of $Y(\fgl(4 \mid 20))$ has
\begin{enumerate}[label=\textup{(\roman*)}]
\item matrix size $(\dim V)^2 \times (\dim V)^2 = 576 \times 576$;
\item a single pole at $u = 0$ of order $1$;
\item eigenvalue factorisation
  $\det R(u) = (1+\hbar u^{-1})^{280}(1-\hbar u^{-1})^{296}$
  on $V \otimes V$, corresponding to the
  $(\mathrm{Sym}^2_s, \Lambda^2_s)$-eigenspace decomposition of
  dimensions $(280, 296)$;
\item after supertrace gauge normalisation (Gow~2007, Prop.~3.3),
  net spectral-parameter polynomial degree $(a, b) = (4, 20)$ on
  each tensor factor with $a+b = \dim V = 24$ and $a-b =
  \mathrm{sdim}(V) = -16$.
\end{enumerate}
\end{lemma}

\begin{proof}
The super-permutation $P_s$ has eigenvalues $\pm 1$ on
$V \otimes V$, with $+1$-eigenspace the super-symmetric square
$\mathrm{Sym}^2_s(V)$ of dimension
$\binom{m+1}{2}+mn+\binom{n}{2} = 10+80+190 = 280$, and
$-1$-eigenspace the super-antisymmetric square
$\Lambda^2_s(V)$ of dimension
$\binom{m}{2}+mn+\binom{n+1}{2} = 6+80+210 = 296$,
summing to $(m+n)^2 = 576$. The super-$R$-matrix
$R(u) = \Id + \hbar u^{-1} P_s$ acts as
$(1+\hbar u^{-1}) \cdot \Id$ on $\mathrm{Sym}^2_s(V)$ and
$(1-\hbar u^{-1})\cdot \Id$ on $\Lambda^2_s(V)$, with
$\det R(u) = (1+\hbar u^{-1})^{280}(1-\hbar u^{-1})^{296}$. After
the standard normalisation of Gow~2007 Prop.~3.3 (dividing by
$\det$ and taking the supertrace-preserving gauge), the residual
spectral-parameter polynomial degree on each of the two tensor
factors is $(a, b)$ with $a+b=\dim V = 24$ and $a-b =
\mathrm{sdim}(V) = -16$, giving $a = 4$, $b = 20$.
\end{proof}

\paragraph{$\osp$-descent.}
Under $\theta$-folding, the $\osp$-$R$-matrix gains an additional
pole at $u = -\hbar\kappa_\osp/2 = 9\hbar$ from the $Q$-projector
term. The $\osp$-$R$-matrix spectral degree is thus $(4, 20, 1)$
on the $\pm$-eigenspaces plus the $Q$-eigenspace; total degree
changes from $24$ to $24 + 1 = 25$ in $u^{\pm 1}$, reflecting the
$\osp$-crossing symmetry of the extended $R$-matrix.

\subsection*{Summary and consistency with Vol III cache}

\paragraph{Consistency checks (from CLAUDE.md, top-15).}
\begin{enumerate}[label=\textup{(\roman*)}]
\item $\CoHA(\C^3) = Y^+(\widehat{\fgl}_1)$ is the positive half of
  the affine $\fgl_1$ Yangian; the K3 analogue is
  Conjecture~\ref{conj:k3-coha-theta-fixed},
  $\CoHA(K3) = Y^+_\osp(4 \mid 20) = Y^+(\fgl(4 \mid 20))^{\theta}$,
  with the $\theta$-fixing essential (cache entry 2, with
  $\C^3 \to K3$ generalisation).
\item $\kappa_{\mathrm{cat}}(K3) = 2$ (K3 fibre, Künneth-multiplicative
  base); the super-Yangian lives on the K3 factor and does not
  enter the $\kappa_{\mathrm{cat}}(K3 \times E) = 0$ computation
  directly.
\item The $(4 \mid 20)$ super-grading is topologically forced
  (Proposition~\ref{prop:4-20-forced-by-hodge}); it is not a
  choice of polarisation. Alternative $(m \mid n)$ with $m+n=24$
  do not arise from the Mukai Hodge involution.
\item The super-$R$-matrix of size $576$
  (Lemma~\ref{lem:super-R-degree}) has spectral-parameter degree
  $(4, 20)$ on each tensor factor; the total matrix rank is
  $24 = 4 + 20$, not $576$.
\item Six routes to $G(K3 \times E)$ are different constructions;
  the super-Yangian construction here is the Yangian face
  (\texttt{k3\_yangian\_chapter.tex}, \S~\ref{subsec:k3-super-yangian})
  specialised to the Mukai super-lattice. It is not the CoHA face
  (that is Conjecture~\ref{conj:k3-coha-theta-fixed}); the two faces
  agree at the level of characters by Schiffmann--Vasserot analogy
  and at the level of associative super-algebras only
  conjecturally.
\end{enumerate}

\paragraph{Status tags.}
\begin{itemize}
\item Definition~\ref{def:super-yangian-gl420-rtt} and the
  super-$R$-matrix \eqref{eq:super-R-matrix}: \emph{proved}
  (Nazarov~1991, Gow~2007, applied at $(m,n)=(4,20)$).
\item Lemma~\ref{lem:super-ybe}, Lemma~\ref{lem:super-R-degree},
  Proposition~\ref{prop:4-20-forced-by-hodge}: \emph{proved here}
  (direct computation).
\item Conjecture~\ref{conj:k3-coha-theta-fixed}: \emph{conjectural};
  depends on Schiffmann--Vasserot analogy (for characters) and CY-C
  six-route convergence (for the algebra-level identification); both
  open.
\item $\osp$-descent (end of each cycle): \emph{proved} at the
  level of the $\theta$-involution existence (Molev~2007,
  Arnaudon--Cramp\'e--Doikou--Frappat--Ragoucy~2003); \emph{conjectural}
  for the global identification $\CoHA(K3) = Y^+_\osp(4 \mid 20)$
  (Conjecture~\ref{conj:k3-coha-theta-fixed}, (iii)).
\end{itemize}

\paragraph{Open questions.}
\begin{enumerate}
\item Does the Chevalley involution $\theta$ on
  $Y(\fgl(4 \mid 20))$ induce an involutive Hopf super-algebra
  automorphism with fixed subalgebra isomorphic to the
  ACDFR reflection-RTT super-algebra $Y_\osp(4 \mid 20)$?
\item Does the super-shuffle product on $Y^+(\fgl(4 \mid 20))$
  restrict to the Schiffmann--Vasserot shuffle product on
  $\CoHA(K3)$ under the $\theta$-fixed embedding?
\item Is the MO $R$-matrix on $\Hilb^n(K3)$ the restriction of a
  universal $R$-matrix of $Y(\fgl(4 \mid 20))$ to a specific
  representation, and if so, is it the $\theta$-fixed evaluation
  or a full-super representation?
\item What is the relationship between $Y(\fgl(4 \mid 20))$ and the
  BKM Borcherds algebra $\fg_{\Delta_5}$ of Gritsenko--Nikulin at
  $N = 1$? The candidate bridge is via the Drinfeld double
  $\mathbf D_\hbar(\fg_{\Delta_5})$ and its Mukai-finite subalgebra.
\end{enumerate}

\subsection*{Primary-literature anchors}

\begin{itemize}
\item Nazarov, M.~L., \emph{Quantum Berezinian and the classical
  Capelli identity}, Lett.\ Math.\ Phys.~21 (1991), 123--131.
\item Nazarov, M.~L., \emph{A mixed hook-length formula for affine
  Hecke algebras}, European J.\ Combin.~23 (2002), 453--466.
\item Gow, L., \emph{Gauss decomposition of the Yangian
  $Y(\fgl_{m|n})$}, Comm.\ Math.\ Phys.~276 (2007), 799--825.
\item Molev, A., \emph{Yangians and Classical Lie Algebras},
  Mathematical Surveys and Monographs 143, AMS, Providence, RI,
  2007. Chapter~3 (super-Yangians and centres).
\item Arnaudon, D., Cramp\'e, N., Doikou, A., Frappat, L., Ragoucy,
  E., \emph{Classification of reflection matrices related to
  $\mathrm{osp}(m \mid 2n)$}, Nucl.\ Phys.\ B~668 (2003), 469--505.
\item Molev, A., Ragoucy, E., \emph{The Yangian
  $Y_\osp(m \mid 2n)$ at the centre}, Algebras and Representation
  Theory~12 (2009), 187--218.
\item Yamane, H., \emph{On defining relations of affine Lie
  superalgebras and affine quantized universal enveloping
  superalgebras}, Publ.\ RIMS Kyoto~35 (1999), 321--390.
\item Stukopin, V., \emph{Yangians of classical Lie superalgebras:
  basic constructions, quantum double and universal $R$-matrix},
  Proc.\ Inst.\ Math.\ NAS Ukraine, 2 (2013), 1195--1201.
\item Schiffmann, O., Vasserot, E., \emph{The elliptic Hall algebra,
  cherednik Hecke algebras and Macdonald polynomials},
  Compos.\ Math.~147 (2011), 188--234; and \emph{Cherednik algebras,
  W-algebras and the equivariant cohomology of the moduli space of
  instantons on $\mathbb A^2$}, Publ.\ Math.\ IH\'ES~118 (2013), 213--342.
\item Costello, K., Yamazaki, M., \emph{Gauge theory and integrability
  III}, arXiv:1908.02289.
\end{itemize}
